\section{Theoretische Grundlagen}

\subsection{Informationssysteme in Unternehmen}
Informationssysteme (IS) sind nach Gabriel im engeren Sinne computergestütze Anwendungssysteme, die von Menschen und Maschinen (beispielsweise Anwendungen) erzeugte Informationen erzeugen und diese miteinander in Beziehung bringen. IS, sowie auch Informations- und Kommunikationssysteme berücksichtigen hierbei die Erfassung, Verarbeitung, Speicherung und Übertragung von Informationen.\autocite[Vgl.][]{InformationssystemEnzyklopaedieWirtschaftsinformatik2019}
Innerhalb von Unternehmen können (betriebliche) Informationssysteme die Prozesse und den Informationsaustausch innerhalb der Unternehmensbereiche, sowie nach außen unterstützen. Die Kosten für zum Speichern von Informationen in Form von Daten auf Hardware sind über die letzten Jahrzehnte kontinuierlich gesunken, während die Computer und Anwendungen leistungsfähiger und komplexer wurden. Diese Entwicklung sorgt für die Möglichkeit, auch kleine Informatiosentitäten zu speichern und zu verarbeiten beziehungsweise in Kontent mit anderen Informationen maschinell auszuwerten.\autocite[Vgl.][Kap. 1.1]{WirtschaftsinformatikEBSCO}

Die Entwicklung, dass Informationen in immer größerem Maßstab gespeichert werden, birgt jedoch auch die Herausforderung, dass eine Suche nach spezifischen Informationen deutlich komplexer und, je nach Systemgröße, oft kaum noch manuell durchführbar ist. Dadurch hat die Rolle der Informationsbeschaffung innerhalb des IS eine höhere Relevanz gewonnen, welche computergestützt, beispielsweise durch Kategorisierung durch Tags, Suche nach Stichwörtern oder auch durch sprachnahe Suche über Künstliche Intelligenz(//Rewrite). Letzteres gewinnt durch die aktuelle Entwicklung im Bereich der KI derzeit stetig an Relevanz. //TODO Quelle

\subsubsection{Implementierungsansätze von Informationssystemen}
Obwohl die Implementierung von Informationssystemen insbesondere den technischen Aspekt fokusiert, sind auch die systemischen, strategischen und oraganisatorischen Aspekte relevant.\autocite[Vgl.][]{wainwrightThreeDomainsImplementing2004}


\subsubsection{Forschungsstand zur Nutzung von LLM in Informationssystemen}
Todo

\subsection{Grundlagen generativer KI-Systeme}
Todo

\subsubsection{Funktionsweise von generativer KI und LLM}
Todo

\subsubsection{Funktionsweise der Informationsverarbeitung und Antworterzeugung in KI-Systemen}
Todo

\subsubsection{Risiken durch implizite Wissensaggregation (Inference Leakage) }
Todo

\subsection{Berechtigungs- und Zugriffskonzepte}
Todo

\subsubsection{Rollenbasierte Zugriffskontrolle}
Todo

\subsubsection{Attributbasierte Zugriffskontrolle }
Todo

\subsubsection{Hybride Berechtigungsmodelle}
Todo

\subsection{Active Directory \& IAM}
Todo

\subsubsection{Active Directory als zentrales IAM-System}
Todo

\subsubsection{Herausforderungen historisch gewachsener Berechtigungsstrukturen}
Todo

